\documentclass[../../../include/open-logic-section]{subfiles}

\begin{document}

\olfileid{sth}{ordinals}{replacement}
\olsection{ప్రతిస్థాపన}

\olref[sth][ordinals][vn]{sec}లో,
క్రమసంఖ్యలను సుక్రమాలకు
నియత ప్రతినిధులైన
క్రమరకాలుగా పరిగణించవచ్చని
సూచిస్తూ వాటిని పరిచయం
చేశాం. అది పని చేయాలంటే
\emph{ప్రతి సుక్రమమూ ఏదో
క్రమసంఖ్యతో సమరూపమని}
నిరూపించాలి. అప్పుడు
$\tuple{A, <} \isomorphic \alpha$
అయ్యే క్రమసంఖ్య $\alpha$ను
$\ordtype{A, <}$గా
నిర్వచించవచ్చు.

దురదృష్టవశాత్తూ,
ఇప్పటివరకు పరిచయం
చేసిన స్వీకృతాలతో
మాత్రమే కావలసిన
ఫలితాన్ని \emph{నిరూపించలేం}.
(ఎందుకో \olref[replacement][strength]{sec}లో
చూస్తాం; ప్రస్తుతానికి
ఇది సాధ్యం కాదనేది
ముఖ్య విషయం.) కొత్త
ఆలోచన అవసరం; అది ఇది:

\begin{axiom}[ప్రతిస్థాపన స్వీకృత పథకం]
ఏ $\phi(x, y)$ సూత్రానికైనా
ఈ కింది వాక్యం ఒక స్వీకృతం:
\begin{quote}
	ఏ $A$కైనా,
	$(\forall x \in A)\lexists![y][\phi(x,y)]$
	అయితే,
	$\Setabs{y}{(\exists x \in A)\phi(x,y)}$
	ఉంటుంది.
\end{quote}
\end{axiom}
\noindent
వేరుచేయడం మాదిరిగానే,
ఇదీ పథకం: భిన్నమైన
$\phi$ సూత్రాలు అనంతంగా
ఉన్నాయి కాబట్టి ఇది
అనంతంగా అనేక స్వీకృతాలను
ఇస్తుంది. దీనిని
సమానార్థకంగా ఈ కింది
విధంగా కూడా రాయవచ్చు;
సాధారణంగా అలాగే రాస్తారు:

\begin{defish}
``$B$'' లేని ఏ
$\phi(x,y)$ సూత్రానికైనా
ఈ కింది వాక్యం ఒక స్వీకృతం:
\[
\forall A[(\forall x \in A)\lexists![y][\phi(x,y)] \lif \exists B\forall y (y \in B \liff (\exists x \in A)\phi(x,y))]
\]
\end{defish}

మొదటిసారి చూసినప్పుడు
ఈ సూత్రం కొంత చిక్కుగా
కనిపిస్తుంది. ప్రతిస్థాపనకు
ఈ కింది తక్షణ పర్యవసానం
మనం అనుసరిస్తున్న
సహజ భావాన్ని బహుశా
\emph{మరింత స్పష్టంగా}
వ్యక్తపరుస్తుంది:\footnote{పదం
అంటే (దాని స్వేచ్ఛా
చరాలకు ఏ విలువలు
ఇచ్చినా) ఖచ్చితంగా
ఒక వస్తువును సూచించే
వ్యక్తీకరణ. ఉదాహరణకు,
``$\emptyset$'' శూన్య
సమితిని సూచించే పదం;
``$\{x\}$'', $x$ ఏదైనా,
$x$ యొక్క ఏక-మూలక సమితిని
సూచించే పదం; ``$x \cup y$'',
$x$, $y$ ఏవైనా,
వాటి సమ్మేళనాన్ని
సూచించే పదం.}

\begin{cor}
ఏ పదం $\tau(x)$,
ఏ సమితి $A$కైనా
ఈ సమితి ఉంటుంది:
\[
	\Setabs{\tau(x)}{x \in A} = \Setabs{y}{(\exists x \in A)y = \tau(x)}.
\]
\end{cor}

\begin{proof}
$\tau$ ఒక \emph{పదం}
కాబట్టి,
$\forall x \lexists![y][\tau(x) = y]$.
ముఖ్యంగా,
$(\forall x \in A)\lexists![y][\tau(x) = y]$.
అందువల్ల ప్రతిస్థాపన
ప్రకారం
$\Setabs{y}{(\exists x \in A)\tau(x) = y}$
ఉంటుంది.
\end{proof}
\noindent
దీని వల్ల ``ప్రతిస్థాపన''
స్వీకృతానికి తగిన పేరని
తెలుస్తుంది: $A$ సమితి
ఇచ్చినప్పుడు, $A$లోని ప్రతి
మూలకాన్ని $\tau$ కింద
దాని చిత్రంతో మార్చి
$\Setabs{\tau(x)}{x \in A}$
అనే కొత్త సమితిని
ఏర్పరచవచ్చు. నిజానికి
ప్రమేయం కింద సమితి
చిత్రపు సంకేతలిపిని
అనుసరించి,
$\Setabs{\tau(x)}{x \in A}$ను
$\funimage{\tau}{A}$గా
రాయవచ్చు.

అయితే ముఖ్యంగా $\tau$
ఒక \emph{పదం}.
\olref[sfr][fun][rel]{sec}
అర్థంలో---అంటే క్రమయుగ్మాల
ఒక నిర్దిష్ట సమితి
అర్థంలో---అది \emph{ప్రమేయానికి}
పేరుగా ఉండవలసిన అవసరం
లేదు. ఆ అర్థంలో
$f$ ఒక ప్రమేయం అయితే,
$\funimage{f}{A} = \Setabs{f(x)}{x \in
A}$ కేవలం $\ran{f}$
యొక్క ఒక ఉపసమితి;
అది ఉంటుందని $\Zminus$
స్వీకృతాలే ఇప్పటికే
నిర్ధారిస్తాయి.\footnote{$\Setabs{y \in \bigcup \bigcup f}{(\exists x \in A)y =
f(x)}$ను పరిశీలించండి.}
దీనికి భిన్నంగా,
ప్రతిస్థాపన మన
స్వీకృతాలకు ఒక
\emph{శక్తివంతమైన}
చేర్పు. దీనిని
\olref[sth][replacement][]{chap}లో
చూస్తాం.


\end{document}
